\chapter*{Introduction générale}
\addcontentsline{toc}{chapter}{Introduction générale}

La transformation numérique des restaurants ne se limite plus à la présence en ligne ou à la prise de réservation. Elle concerne désormais l'organisation opérationnelle dans son ensemble : coordination entre la salle et la cuisine, disponibilité des plats, suivi des stocks, gestion des paiements, fidélisation et analyse des retours clients. Dans un établissement où plusieurs acteurs interviennent simultanément, une information mal transmise peut entraîner un retard, une erreur de préparation, une rupture non anticipée ou une dégradation de l'expérience client.

Le projet Tastify s'inscrit dans ce contexte. Il propose une application web centralisée pour la gestion d'un restaurant, en couvrant à la fois les besoins internes du personnel et les interactions avec les clients. La solution s'appuie sur une architecture full-stack composée d'un backend Django REST Framework, d'une base de données MySQL, de Redis et Celery pour les échanges temps réel et les traitements asynchrones, de Django Channels pour les WebSocket, ainsi que de deux applications React : une interface staff et un portail client.

La particularité du projet réside dans son périmètre transversal. Tastify ne se limite pas à un module de menu ou à une caisse numérique : il relie les tables, les commandes, le Kitchen Display System, les stocks, les réservations, les paiements, la fidélité, les avis clients et les tableaux de bord. Cette intégration rend le projet pertinent dans le cadre d'un Projet de Fin d'Année, car elle mobilise plusieurs dimensions du génie logiciel : modélisation, architecture, API, sécurité, temps réel, asynchronisme, expérience utilisateur et validation.

Le présent rapport est organisé selon une progression académique classique. Le premier chapitre présente le contexte général, la problématique, l'étude de l'existant et les objectifs. Le deuxième chapitre formalise l'analyse des besoins et la conception. Un chapitre spécifique est consacré à l'analyse de sentiment, car cette fonctionnalité combine un enjeu métier et un pipeline technique précis. Les chapitres suivants décrivent l'architecture, la réalisation, les interfaces, les tests, les limites et les perspectives.

La rédaction adopte enfin un principe de rigueur : les fonctionnalités sont présentées lorsqu'elles sont effectivement intégrées au périmètre du projet, tandis que les limites et les éléments restant à compléter sont explicitement distingués. Cette approche permet de défendre le travail réalisé sans surestimer son niveau de maturité.
